\scanpage{79}
de \unsure{passgge|passsge|passage} au quotient (sauf erreur de signe, car je n'ai nulle part fait attention aux signes).

8.4. Considérons maintenant l'homomorphisme
\[
\handwritten{\underline}{G}(\handwritten{\overline}{\unsure{\handwritten{\underline}{F}|\handwritten{\underline}{\mathbb{F}}}}_p)\handwritten{\longrightarrow}\handwritten{\underline}{G}(\handwritten{\underline}{F}_p)\quad .
\tag*{(8.4.1)}
\]
Il est clair qu'il se factorise par $\handwritten{\underline}{H}$ (i.e. que le composé
\[
\handwritten{\underline}{Q}_p\text{-faisceaux constructibles sur }\handwritten{\underline}{F}_p\ \overset{i}{\handwritten{\longrightarrow}}\ \operatorname{Fcriso}(\handwritten{\underline}{F}_p)\handwritten{\longrightarrow}\operatorname{Fcriso}(\handwritten{\overline}{\handwritten{\underline}{F}}_p)
\]
s'envoie dans les objets ``constants'' ...) , d'où un homomorphisme
\[
\handwritten{\underline}{G}(\handwritten{\overline}{\handwritten{\underline}{F}}_p)\handwritten{\longrightarrow}\handwritten{\underline}{H}\quad .
\tag*{(8.4.2)}
\]
Par Dieudonné-Manin, on connait la structure des objets de $\operatorname{Fcriso}(\handwritten{\overline}{k})$ pour $\handwritten{\overline}{k}$ algébriquement clos\deleted{\unsure{x|e}}, on sait que ces objets sont semi-simples et proviennent (mod isomorphisme) d'objets de $\operatorname{Fcriso}(\handwritten{\underline}{F}_p)$. Cela implique que (8.4.1), donc aassi (8.4.2), est un monomorphisme de gerbes. Pour montrer que c'est une équivalence, il suffit de montrer qu'un objet simple non constant de \deleted{$\handwritten{\underline}{H}$ me donne} $\operatorname{Rep}(\handwritten{\underline}{H})$ donne un objet non constant de $\operatorname{Rep}(\handwritten{\underline}{G}(\handwritten{\overline}{k}))\handwritten{\overset{\approx}{\longrightarrow}}\operatorname{Fcriso}(\handwritten{\overline}{k}))$. Or les \unsure{cla sees|cla ses|classes} d'objetssimples de $\operatorname{Rep}(\handwritten{\underline}{H})$ sont indexés par les éléments de $\handwritten{\underline}{Q}$, et si on considère l'objet simple de $\operatorname{Fcriso}(\handwritten{\underline}{F}_p)\handwritten{\simeq}\operatorname{Rep}(\handwritten{\underline}{G}(\handwritten{\underline}{F}_p^{\deleted{\handwritten{\unsure{*|\overline{\vphantom{F}}}}}}))$ défini par $(E,\handwritten{\underline}{F}_E)$, avec $\handwritten{\underline}{F}_E$ ayant comme polynôme minimal $T^r-p^s$ avec $r,s\handwritten{\in}\handwritten{\underline}{Z}$, $r\geq1$, $(r,s)=1$, alors sa restriction à $\handwritten{\underline}{H}$ donne un objet simple \deleted{\unsure{de|du}} (car il reste simple dans $\operatorname{Fcriso}(\handwritten{\overline}{k})$) de classe $s/r\handwritten{\in}\handwritten{\underline}{Q}$, et celui\handwritten{-}ci, d'après la classification de Dieudonné-Manin, \deleted{n'est pas} a une image dans $\operatorname{Fc\ iso}(\handwritten{\overline}{k})$ qui n'est pas constante si $r\handwritten{\ne}0$ ...

\handwritten{8.5.} On a donc \deleted{déterminé} résolu par l'affirmative, pour $k=\handwritten{\underline}{F}_p$, la question 6.2, en déterminant en même temps (à équivalence définie à isomorphisme unique près) la gerbe $\unsure{\handwritten{\underline}{G}|\handwritten{\underline}{H}}(\handwritten{\overline}{k})$ associée à une clôture algébrique $\handwritten{\overline}{k}$ de $k=\handwritten{\underline}{F}_p$ ; \deleted{comme pour} elle ne dépend pas (à \deleted{\unsure{isomorphi|isomorphix}} équivalence déterminée à isomorphisme unique près) du choix de la clôture algébrique envisagée $\handwritten{\overline}{k}$, et par Dieudonné-Manin est can. équi%
% archive sheet number: 78
